% Pursuing Stacks (cotes 134-2 to 134-8): archive cote 134-6, site page 37 (PDF page 38), typescript page 355.
% typescript. Draft compared: ps5.tex:3153-3221.
% AI-transcribed by Claude Opus 5.5 (claude-opus-5-5, adaptive thinking, effort high), Message Batches API, batch ps3 (msgbatch_01XdFmuVVBAcgc7EFxdtCqsj), 2026-11-12; reviewed against both supplied scan halves; see CHECK_DECISIONS.tsv.
\scanpage{134-6/37}
\mspage{355}

The propert\over[typed]{\illegible{1}}{y} \xout{that} for a topos X, with category of sheaves \hand{$\mathcal{A}$}, that
any object F in \hand{$\mathcal{A}$} \xout{such that} covering the final`ob\over[typed]{h}{j}ect e\hand{$_{\mathcal{A}}$} \inserted{should} ha\over{s}{ve} a section,
can be expressed by saying that the latter is a \uline{projective} object in
the category \hand{$\mathcal{A}$}. Following the principle to use the same names for
properties of a topos, and corresponding properties of the final
sheaf on it, we may call a topos with the above property a \uline{projective}
\uline{topos}. Thus the``non-empty'' topo$\overset{\hand{\illegible{1}}}{\text{\i}}$ such that every ``non-empty'' sheaf \xout{\illegible{2}}
has a section, are exactly the topo\"{\i} which are both vertical and projec%
-tive.

\begin{center}\shortstack{*\qquad *\\{}*}\end{center}

Here is the promised ``exact relationship'' between the three
standard homotopy structures h,h',h'' on A\^{}, where A is any small cate-%
gory (cf. end of section 96, p. 354) :

% Hand-drawn arrows are marked blue; typed objects remain plain.
\diagram{\begin{tikzpicture}[x=1cm,y=1cm]
\node (h0) at (0,0) {\ensuremath{h}};
\node (h1) at (3.3,0) {\ensuremath{h'}};
\node (h2) at (6.6,0) {\ensuremath{h''}};
\node (h3) at (9.9,0) {\ensuremath{h}};
\node at (1.65,0) {\ensuremath{\hand{\leq}}};
\node at (4.95,0) {\ensuremath{\hand{\leq}}};
\node at (8.25,0) {\ensuremath{\hand{\leq}}};
\node[inner sep=2pt] (c1) at (1.65,.65) {\ensuremath{\scriptstyle(A\ \text{tot.}\ \underline{W}\text{-asph.})}};
\node at (4.95,.65) {\ensuremath{\scriptstyle(\underline{W}\hand{\subset}\underline{W}_o)}};
\node at (8.25,.45) {\ensuremath{\scriptstyle\xout{(A\ \text{vertical})}}};
\node at (8.25,.95) {\inserted[typed]{\ensuremath{(A\ \text{vertical})}}};
\node at (-.4,.65) {\xout{\illegible{1}}};
\node at (-1.2,-.35) {\ensuremath{(\hand{*})}};
\node[rotate=90] at (0,-.45) {\ensuremath{\hand{=}}};
\node[rotate=90] at (3.3,-.45) {\ensuremath{\hand{=}}};
\node[rotate=90] at (6.6,-.45) {\ensuremath{\hand{=}}};
\node at (0,-.95) {\ensuremath{h_{A\hat{}\deleted{\illegible{1}}\hand{,}A}}};
\node at (3.3,-.95) {\ensuremath{h_{\underline{W}_A}}};
\node at (6.6,-.95) {\ensuremath{h_{\underline{W}_{oA}}=h_{A\hat{}}}};
\node at (9.9,-.95) {\ensuremath{,}};
\node[inner sep=3pt] (ec) at (-2,1.05) {\hand{\'ecarter}};
\draw[color=blue] (ec.center) ellipse (1.05cm and .48cm);
\draw[color=blue] ([xshift=1.05cm]ec.center) .. controls (0,1.25) and (.05,1.32) .. ([xshift=-.1cm,yshift=.1cm]c1.north west);
\draw[color=blue] ([xshift=-.08cm]c1.south west)--([xshift=-.08cm]c1.north west);
\draw[color=blue] ([xshift=.08cm]c1.south east)--([xshift=.08cm]c1.north east);
\draw[color=blue] ([xshift=-.08cm]c1.north west) .. controls ([yshift=.35cm]c1.north) .. ([xshift=.08cm]c1.north east);
\end{tikzpicture}}

where \uline{W} is a given basic localizer. Above each one of the three \inserted[typed]{conditional} ine-%
qualities between homotopy structures h,h',h'' I wrote the natural \xout{assu}
assumption \inserted[typed]{on A or \uline{W}} validating it , and \xout{underneath of} in the diagram I have
recalled the definition of the three homotopy structures. A propos
the descr\over[typed]{n}{i}ption $h = h_{A\hat{},A}$ , the notation used here is $h_{M,A}$ when
M is a category stable under finite products and A a full subcategory,
for designating the homotopy structure on M generated by intervals in
M coming from A. A propos $h'' = h_{A\hat{}}$ , I recall the notation $h_M$ for
designating the canonical homotopy structure on a category M satisfying
suitable conditions (section 57). Also I recall
\[ \underline{W}_o = \hand{\{} f\hand{\in}\mathrm{Fl}(\mathrm{Cat}) \hand{|} \hand{\pi}_o(f) \text{ bijective}\hand{\}} \ . \]
The diagram implies that if $\underline{W}\hand{\subset}\underline{W}_o$ i.e. \uline{W} satisfies Loc 4), and if
moreover A is vertical and \xout{totally $\underline{W}$ non empty (hence} tota\over{\illegible{1}}{l}\over{l}{l}y $\underline{W}_{\xout{oo}}$-as-%
pheric, \inserted[typed]{then all three homotopy structures coincide. Also, taking} \deleted{\xout{by what we just saw, and a fortiori $\underline{W}$-aspheric. If we admit her\unsure{e}}}
\deleted{\xout{that necessarily \unsure{$\underline{W}$ $\underline{W}_{oo}$}}} $\underline{W} = \underline{W}_o$ , we see that h=h'' if A is verti-%
cal and totally 0-connected, w\over{h}{h}\unsure{t|i}ch is \over{\unsure{a}}{l}emma 2 of p.268 for A\^{}, A .

From (1) it follows of course that if A is a contractor, then
(for any \uline{W} satisfying Loc 4)) the three homotopy structures h,h',h''
on A\^{} coincide. In case the contractor A is not trivial, hence A is

\margin{37.}

% Visual collation: both supplied image halves inspected; see CHECK_DECISIONS_134_7_8.tsv.
